\documentclass[11pt]{article}
% ------------------------------------------------------------------
% RESEARCH LETTER version, condensed for Quantitative Finance
% (Taylor & Francis). Stand-in class = article; at submission port to
% the T&F "Interact" template (interact.cls) with the QF reference
% style. Sectioning, abstract and keywords map directly. All numbers
% are the same committed results as the full paper.
% ------------------------------------------------------------------
\usepackage[margin=1in]{geometry}
\usepackage{amsmath,amssymb}
\usepackage{graphicx}
\graphicspath{{../figures/}}
\usepackage{booktabs}
\usepackage[round]{natbib}
\usepackage{xcolor}
\usepackage[hidelinks]{hyperref}
\newcommand{\todo}[1]{\textcolor{red}{[#1]}}
\newcommand{\bp}{\ensuremath{\,\mathrm{bp}}}

\title{\textbf{Meshfree pricing and fast calibration of rough volatility:\\
a physics-informed neural network for the Markovian-lifted PDE}}
\author{\todo{Authors / affiliations}}
\date{Research Letter. \today}

\begin{document}
\maketitle

\begin{abstract}
Rough-volatility models are non-Markovian and admit no finite-dimensional pricing
partial differential equation (PDE); the Markovian lift restores one by
approximating the fractional kernel with $n$ exponential factors, at the cost of
an $(n{+}1)$-dimensional problem. We solve this lifted PDE with a physics-informed
neural network (PINN) and ask how the network's accuracy scales with the lifted
dimension. Validating against an independent lifted Monte Carlo, and separating
the network's \emph{solve-error} from the model's \emph{lift-error} in price
space, we find the solve-error grows only sub-linearly---staying under about
eleven price basis points to $n{=}32$---while a naive finite-difference grid
($50^{n+1}$ points) is already infeasible by $n{=}8$. The crossover between the
two error sources gives a practical rule for choosing $n$. Calibrated to Deribit
Bitcoin options the model returns a rough Hurst exponent ($H\approx0.09$) with
strong negative skew. A parametric PINN that takes the model parameters as inputs
(with the spot variance and maturity fixed) is too coarse to calibrate on its own,
but as a warm start for a short local optimization against the true pricer it
recovers from-scratch fit quality at roughly nine times fewer true-model
evaluations. The message is that where grid methods succumb to dimension, a
mesh-free PINN degrades gracefully, and is most useful in practice as a
differentiable warm start for calibration.
\end{abstract}

\noindent\textbf{Keywords:} physics-informed neural networks; rough volatility;
Markovian lift; curse of dimensionality; option pricing; calibration.

% ==================================================================
\section{Introduction}
Pricing under a stochastic model is, underneath, a PDE, and for the standard
models it is low-dimensional and easily gridded. Rough volatility breaks this:
motivated by the empirical roughness of realized variance and the steep
short-maturity skew \citep{gatheral2018volatility,bennedsen2022decoupling}, the
variance is driven by a fractional noise with Hurst exponent $H<1/2$, which is
non-Markovian and leaves no finite-dimensional pricing PDE
\citep{eleuch2019characteristic}. The Markovian lift approximates the fractional
kernel by a sum of $n$ exponentials and recovers an $(n{+}1)$-dimensional PDE
\citep{abijaber2019multifactor,abijaber2019lifting}, trading intractability for
high dimension---with the approximation order $n$ sitting directly on the PDE
dimension.

That coupling makes the lifted problem an ideal test for mesh-free neural PDE
solvers. Physics-informed neural networks minimize the PDE residual at sampled
points \citep{raissi2019physics} and, like other deep PDE solvers
\citep{sirignano2018dgm,han2018solving}, avoid the exponential cost of a grid; they
also learn from the equation itself (no labelled prices) and are differentiable in
their inputs, which we exploit for calibration. The closest prior work,
\citet{papapantoleon2025timestepping}, prices the lifted (rough) Heston PDE with a
time-stepping deep gradient flow; we are not the first mesh-free solver here. Our
contribution is a PINN residual formulation together with (i)~a dimension-scaling
result---the network's solve-error grows sub-linearly while a naive grid explodes;
(ii)~a solve/lift error decomposition and a crossover rule for $n$;
(iii)~calibration to real Bitcoin options; and (iv)~a parametric-PINN warm-start
that accelerates recalibration, reported honestly. We test everything on Bitcoin
options, a liquid, rough, high-skew market.

% ==================================================================
\section{Method}
Under rough Heston the variance is driven by the kernel
$K(t)=t^{\alpha-1}/\Gamma(\alpha)$, $\alpha=H+\tfrac12$; the lift replaces it with
$K(t)\approx\sum_i c_i e^{-x_i t}$, giving $n$ Ornstein--Uhlenbeck factors $U^i$
and $V=V_0+\sum_i c_i U^i$. The discounted price $P(t,x,u)$, $x=\log S$, solves
\begin{align}
0 = {}& P_t + \big(r-\tfrac12 V\big)P_x + \tfrac12 V P_{xx}
   + \textstyle\sum_i[-x_i u_i + \lambda(\theta-V)]P_{u_i}\nonumber\\
   &+ \tfrac12\nu^2 V\textstyle\sum_{i,j}P_{u_i u_j} + \rho\nu V\textstyle\sum_i P_{xu_i} - rP .
\label{eq:pde}
\end{align}
We solve \eqref{eq:pde} with a $\tanh$ PINN of $(\tau,x,u)$, $\tau=T-t$, using the
ansatz $P=P_{\mathrm{base}}(\tau,x)+\tau N_\theta(\tau,x,u)$ that enforces the
initial condition exactly, with $P_{\mathrm{base}}$ a Black--Scholes price anchored
at the expected integrated variance. We add asymptotic boundary losses and,
because PINNs are prone to optimization pathologies
\citep{krishnapriyan2021characterizing,wang2022when}, self-adaptive per-term loss
weights \citep{kendall2018multitask,mcclenny2023self} and a causal
$\tau$-curriculum \citep{wang2024respecting}; factor collocation is placed by a
short pre-flight Monte Carlo.

We hold the network to two independent ground truths: a semi-analytic Fourier
pricer \citep{eleuch2019characteristic,carr1999option}, exact for the model
($n\to\infty$), and a lifted Monte Carlo (an exponential-integrator Euler scheme),
exact for the finite-$n$ PDE the network actually solves. We therefore never
report a single ``PINN-versus-Fourier'' number but decompose it, per strike and in
price space where the signed errors add exactly,
\begin{equation}
\underbrace{(\mathrm{PINN}{-}\mathrm{Fourier})}_{\text{total}}
=\underbrace{(\mathrm{PINN}{-}\mathrm{MC})}_{\text{solve-error}}
+\underbrace{(\mathrm{MC}{-}\mathrm{Fourier})}_{\text{lift-error}} .
\label{eq:split}
\end{equation}
The solve-error is what training controls; the lift-error is a model approximation
that only larger $n$ can reduce. For calibration we compare differential evolution
with the Fourier pricer, gradient descent on a parametric PINN (model parameters as
inputs, $V_0$ and maturity fixed), and a hybrid that polishes the PINN's guess with
a few Fourier steps. Data are end-of-day BTC option quotes from Deribit (a
cryptocurrency exchange and data provider); we set $r=0$ and infer the forward from
quotes.

% ==================================================================
\section{Results}
\paragraph{The network scales where the grid cannot.}
Training one architecture (width $64$, depth $4$) at $n=4,8,16,32$ on the canonical
BTC calibration ($T=0.15$) over three seeds, the solve-error against the lifted MC
grows only mildly, from $5.4\pm1.0$ to $10.7\pm2.0$ price\bp{}---about a doubling
for an eight-fold rise in dimension (Figure~\ref{fig:nconv}). It is not flat (an
early single-seed run misled us), but grows far more slowly than the dimension;
with three seeds the adjacent-$n$ bars overlap, so the robust statement is the
two-fold growth over the full range. Training cost stays bounded ($6$--$10$ min,
CPU). By contrast a naive full-tensor grid needs $50^{n+1}$ points---$3.1\times
10^{8}$, $2.0\times10^{15}$, $7.6\times10^{28}$, $1.2\times10^{56}$ at
$n=4,8,16,32$---so even $n=8$ ($\sim\!16$ PB) is infeasible. We are fair about this:
$50^{n+1}$ is the \emph{naive} grid, and sparse-grid/ADI/tensor-train schemes reach
further but still degrade with $n$ and are not standard for this PDE, so a
competitive finite-difference solve at $n=16$--$32$ is \emph{impractical}, not
provably impossible.

\begin{figure}[t]\centering
\includegraphics[width=0.74\linewidth]{fig1_n_convergence.png}
\caption{Dimension scaling. PINN solve-error (blue, 3-seed mean\,$\pm$\,SD) grows
sub-linearly and stays under $\sim\!11$ price\bp{} to $n=32$, while the lift-error
(red) falls and training time (green) stays bounded; the naive finite-difference
grid (inset) is infeasible by $n=8$.}
\label{fig:nconv}
\end{figure}

\paragraph{Decomposition and a rule for $n$.}
Split by \eqref{eq:split}, the historical ``PINN vs Fourier'' gap of $114$--$127$
vol\bp{} at $n=4$ is almost entirely lift: $37$ vol\bp{} of solve-error on $133$ of
lift (partly cancelling to $121$), and only $8$ of solve-error with correlation
switched off. The network already solves its PDE to well under $50$ vol\bp. Since
the lift-error shrinks with $n$ ($133,63,27,12$ vol\bp{}) while the solve-error
climbs gently, they cross near $n=8$--$10$ at $T\approx0.15$: below it add factors,
above it stop. The crossover moves with maturity---the lift-error grows with $T$
(e.g.\ $8.6\to31.8$ price\bp{} at $n=4$ as $T:0.05\to0.5$), so longer maturities
need larger $n$.

\paragraph{Bitcoin is rough.}
A full-budget calibration to the Deribit BTC surface (differential evolution,
$60$ generations plus polish; $13$ maturities, $97$ points) returns $H=0.090$,
$\rho=-0.792$, $V_0=0.103$, $\theta=0.253$, $\lambda=2.567$, $\nu=0.300$ at a
vega-weighted RMSE of $496$ vol\bp{}. The Hurst exponent is firmly rough and the
skew strongly negative. The fit is deliberately coarse: at $496$ vol\bp{} it
reflects a known limitation of the lifted rough-Heston form on wide crypto smiles,
a model-expressiveness issue orthogonal to the solve-error.

\paragraph{A parametric network and a warm-start hybrid.}
A single network over the five-parameter box costs accuracy: its solve-error at the
box centre is $11.0\pm4.7$ price\bp{} versus $5.9\pm1.2$ for a point-trained network
at the same $n=10$, and worse in the rough corners (up to $47.7$). We therefore do
not use it as a standalone calibrator. Calibrating the BTC $T\approx0.15$ slice
three ways (Table~\ref{tab:calib}), Fourier-in-the-loop reaches $23.7$ vol\bp{} in
$2018$ Fourier evaluations ($\sim\!168$ s); the PINN alone is fast ($\sim\!2$ s,
zero Fourier calls) but poor ($222\pm20$ vol\bp{} under the true model, partly
because the slice optimum lies outside the training box). Seeding a short Fourier
polish with the PINN's guess reaches $29.4\pm7.1$ vol\bp{}---on par with
from-scratch on a $34$--$52\%$ smile---in only $237\pm46$ Fourier evaluations
($\sim\!23$ s): about nine times fewer true-model calls and seven times faster. The
parametric network is thus most useful as a differentiable warm start.

\begin{table}[t]\centering
\caption{Calibration of the BTC $T\approx0.15$ slice, three ways (3 PINN seeds; all
judged by the true Fourier model over the same feasible region).}
\label{tab:calib}
\begin{tabular}{lccc}
\toprule
method & fit RMSE (vol\bp) & wall-clock (s) & Fourier evals\\
\midrule
Fourier-loop from scratch        & $23.7$       & $168$   & $2018$\\
PINN-only (single-shot)          & $222\pm20$   & $\sim2$ & $0$\\
Hybrid (PINN $+$ Fourier polish) & $29.4\pm7.1$ & $23$    & $237\pm46$\\
\bottomrule
\end{tabular}
\end{table}

% ==================================================================
\section{Conclusion}
On the lifted rough-Heston PDE, a physics-informed neural network's solve-error
grows only sub-linearly in the lifted dimension while naive finite differences
become infeasible---a concrete curse-of-dimensionality advantage on an
economically real problem. Separating solve-error from lift-error is essential to
read the numbers correctly and yields a crossover rule for the factor count.
Applied to real Bitcoin options, the model is genuinely rough, and a parametric
PINN, though not accurate enough to calibrate alone, is an effective warm start
that recalibrates at roughly an order of magnitude fewer true-model evaluations.
Natural next steps are adding the maturity as an input for whole-surface pricing,
larger networks and GPUs to shrink the parametric penalty, and extensions to other
rough models and path-dependent payoffs.

\bibliographystyle{plainnat}
\bibliography{refs}
\end{document}
